\subsection{Nullary Operations}
\label{subsec:nullary-operations}

\begin{definition}[Nullary Operation]
\label{def:nullary-operation}
Let $S$ be a set.
A \textbf{nullary operation} on $S$ is the selection of a distinguished element
of $S$.
\end{definition}

\begin{remark*}[Standard quantified statement]
\[
\operatorname{NullaryOperation}(c,S)
\Longleftrightarrow
c\in S.
\]
\end{remark*}

\begin{remark*}[Definition predicate reading]
$c$ is a nullary operation on $S$ exactly when it selects a distinguished
element of $S$.
\end{remark*}

\begin{remark*}[Negated quantified statement]
\[
\neg\operatorname{NullaryOperation}(c,S)
\Longleftrightarrow
c\notin S.
\]
\end{remark*}

\begin{remark*}[Interpretation]
In model-theoretic language, a constant symbol is interpreted as a nullary
operation.  The identity elements $0$, $1$, and $e$ are typical algebraic
examples.
\end{remark*}

\begin{remark*}[Dependencies]
\begin{itemize}
  \item Carrier set.
\end{itemize}
\end{remark*}
